\documentclass[10pt]{article}
% Resumen para consultar durante el Examen 1 de Óptica No Lineal (temas 1–4).
% Versión de una columna, con explicación de cada tema (2026-11-15).
% Compilar con XeLaTeX (dos pasadas):  xelatex resumen-examen-1.tex
\usepackage[spanish,es-noshorthands,es-nodecimaldot,es-tabla]{babel}
\usepackage{fontspec}
\usepackage{amsmath,amssymb,mathtools}
\usepackage{booktabs,array}
\usepackage{xcolor}
\usepackage{enumitem}
\usepackage[letterpaper,margin=2cm,top=2cm,bottom=2cm]{geometry}
\usepackage{titlesec}
\usepackage{fancyhdr}
\usepackage[most]{tcolorbox}
\usepackage{hyperref}
\hypersetup{colorlinks=true,linkcolor=black,urlcolor=black}

\definecolor{tema}{RGB}{20,60,120}
\definecolor{caja}{RGB}{235,241,250}
\definecolor{idea}{RGB}{245,245,235}
\definecolor{alerta}{RGB}{170,30,30}
\titleformat{\section}{\Large\bfseries\color{tema}}{\thesection.}{0.5em}{}[\vspace{-0.4ex}\color{tema}\rule{\linewidth}{0.8pt}]
\titleformat{\subsection}{\large\bfseries\color{tema}}{\thesubsection}{0.5em}{}
\titlespacing*{\section}{0pt}{3ex}{1.5ex}
\titlespacing*{\subsection}{0pt}{2.2ex}{0.8ex}
\setlist{itemsep=0.3ex,topsep=0.6ex,leftmargin=1.5em}
\setlength{\parindent}{0pt}
\setlength{\parskip}{0.8ex}
\pagestyle{fancy}\fancyhf{}
\lhead{\small Óptica No Lineal 2026-II — Resumen Examen 1 (temas 1–4)}
\rhead{\small p.~\thepage}
\renewcommand{\headrulewidth}{0.4pt}
\setlength{\headheight}{14pt}

\newcommand{\Ec}{\boldsymbol{\mathcal E}}
\newcommand{\Hc}{\boldsymbol{\mathcal H}}
\newcommand{\Dc}{\boldsymbol{\mathcal D}}
\newcommand{\Bc}{\boldsymbol{\mathcal B}}
\newcommand{\Pc}{\boldsymbol{\mathcal P}}
\newcommand{\Sc}{\boldsymbol{\mathcal S}}
\newcommand{\E}{\mathbf E}
\newcommand{\Hv}{\mathbf H}
\newcommand{\D}{\mathbf D}
\newcommand{\kv}{\mathbf k}
\newcommand{\rv}{\mathbf r}
\newcommand{\J}{\mathbf J}
\newcommand{\T}{\mathbf T}
\newcommand{\R}{\mathbf R}
\newcommand{\ud}{\hat{\mathbf u}}
\newcommand{\cc}{\text{c.c.}}
\newcommand{\ojo}[1]{{\color{alerta}\textbf{Ojo:} #1}}

% Caja de fórmulas clave
\newtcolorbox{clave}[1][]{enhanced,breakable,colback=caja,colframe=tema,boxrule=0.5pt,arc=2pt,
  left=6pt,right=6pt,top=4pt,bottom=4pt,fonttitle=\bfseries\small,#1}
% Caja de idea física (qué significa, por qué importa)
\newtcolorbox{idea}{enhanced,breakable,colback=idea,colframe=black!25,boxrule=0.4pt,arc=2pt,
  left=6pt,right=6pt,top=3pt,bottom=3pt,title={Idea física},fonttitle=\bfseries\small,
  coltitle=black,colbacktitle=black!8}

\begin{document}

\begin{center}
{\LARGE\bfseries\color{tema} Óptica No Lineal — Resumen para el Examen 1}\\[1ex]
{\normalsize Temas 1–4: ondas electromagnéticas (Saleh 5.1, 5.2, 5.4, 2.2) · haz gaussiano, absorción y dispersión, pulsos (Saleh 3.1, 5.5, 5.6) · polarización y cristales (Saleh cap.~6) · susceptibilidad no lineal (Boyd cap.~1)}\\[1ex]
{\small Juan E. Rodríguez Ochoa · UNAL · Prof. R. Acuña Herrera\\
Saleh \& Teich, \emph{Fundamentals of Photonics}, 2.ª ed. (2007); Boyd, \emph{Nonlinear Optics}, 3.ª ed. (2008).\\
Los números entre paréntesis, p.~ej. (5.2-25) o (1.4.24), son ecuaciones del libro correspondiente.}
\end{center}

\tableofcontents

%=====================================================================
\section*{0. Constantes y convenciones (leer primero)}
\addcontentsline{toc}{section}{0. Constantes y convenciones}
%=====================================================================

\subsection*{Constantes físicas}
\begin{center}
\renewcommand{\arraystretch}{1.25}
\begin{tabular}{@{}ll@{\qquad}ll@{}}
\toprule
$c_0$ & $2.99792458\times10^{8}$ m/s & $e$ & $1.602\times10^{-19}$ C\\
$\epsilon_0$ & $8.854\times10^{-12}$ F/m & $m_e$ & $9.109\times10^{-31}$ kg\\
$\mu_0$ & $4\pi\times10^{-7}$ H/m & $h$ & $6.626\times10^{-34}$ J·s\\
$\eta_0=\sqrt{\mu_0/\epsilon_0}$ & $376.73\ \Omega\approx120\pi\ \Omega$ & $\hbar$ & $1.0546\times10^{-34}$ J·s\\
$1/(4\pi\epsilon_0)$ & $8.988\times10^{9}$ N·m²/C² & $a_0$ (radio de Bohr) & $0.529$ Å\\
 & & $k_B$ & $1.381\times10^{-23}$ J/K\\
\bottomrule
\end{tabular}
\end{center}
Relaciones básicas: $\omega=2\pi\nu$, \ $k_0=\omega/c_0=2\pi/\lambda_0$; en un medio de índice $n$: $k=nk_0$, $\lambda=\lambda_0/n$, $c=c_0/n$.

\subsection*{Las dos convenciones: Saleh frente a Boyd}
Los temas 1–3 usan el Saleh y el tema 4 usa el Boyd. Los dos libros escriben los campos complejos de forma distinta, así que una fórmula copiada de uno al otro sin convertir sale con un factor 2 o un signo equivocado. La tabla resume las diferencias:

\begin{center}
\renewcommand{\arraystretch}{1.35}
\begin{tabular}{@{}p{0.16\linewidth}p{0.38\linewidth}p{0.38\linewidth}@{}}
\toprule
 & \textbf{Saleh (temas 1–3)} & \textbf{Boyd (tema 4)}\\\midrule
Dependencia temporal & $e^{+j\omega t}$; onda hacia $+z$: $e^{-jkz}$ & $e^{-i\omega t}$; onda hacia $+z$: $e^{+ikz}$\\
Campo real & $\Ec=\mathrm{Re}\{\E e^{j\omega t}\}$ ($\E$ es la amplitud pico) & $\tilde E=Ee^{-i\omega t}+\cc$ ($|E|$ es la mitad del pico)\\
Intensidad & $I=|E|^2/2\eta$ & $I=2n\epsilon_0c|E|^2$\\
Absorción & $\chi''<0\Rightarrow\alpha>0$ & $\mathrm{Im}\,\chi>0\Rightarrow$ absorbe\\
Polarización no lineal & $\mathcal P_{NL}=2\mathrm d\,\mathcal E^2+4\chi^{(3)}_S\mathcal E^3$ & $\tilde P=\epsilon_0[\chi^{(1)}\tilde E+\chi^{(2)}\tilde E^2+\chi^{(3)}\tilde E^3]$\\
\bottomrule
\end{tabular}
\end{center}

\begin{clave}[title={Equivalencias}]
$E_{\rm Saleh}=2E_{\rm Boyd}$; \quad $j\leftrightarrow -i$; \quad $\mathrm d_{\rm Saleh}=\tfrac12\epsilon_0\chi^{(2)}=\epsilon_0d_{\rm Boyd}$; \quad
$\chi^{(3)}_{\rm Saleh}=\tfrac14\epsilon_0\chi^{(3)}_{\rm Boyd}$; \quad $n_2$ [m²/W] es el mismo en ambos libros.
\end{clave}
\ojo{al pasar de un libro al otro cambia el signo de la parte imaginaria de $\chi$ (porque $j\to-i$) y aparece el factor 2 en las amplitudes.}

%=====================================================================
\section{Ondas electromagnéticas (Saleh 5.1–5.4, 2.2)}
%=====================================================================

\begin{idea}
La luz es una onda electromagnética. Todo lo que pasa en un material (refracción, absorción, dispersión, birrefringencia y los efectos no lineales) entra en las ecuaciones de Maxwell por un único término: la \textbf{polarización} $\Pc$, el momento dipolar por unidad de volumen que el campo induce en el medio. Este tema construye la ecuación de onda con $\Pc$ como fuente; los temas siguientes solo cambian la relación entre $\Pc$ y $\Ec$.
\end{idea}

\subsection{Ecuaciones de Maxwell, condiciones de frontera y energía (§5.1)}
En un medio sin cargas ni corrientes libres:
\[
\nabla\times\Hc=\frac{\partial\Dc}{\partial t},\qquad
\nabla\times\Ec=-\frac{\partial\Bc}{\partial t},\qquad
\nabla\cdot\Dc=0,\qquad \nabla\cdot\Bc=0 \qquad(5.1\text{-}7\ \text{a}\ 5.1\text{-}10)
\]
Las relaciones constitutivas dicen cómo responde el medio:
\[
\Dc=\epsilon_0\Ec+\Pc,\qquad \Bc=\mu_0\Hc+\mu_0\boldsymbol{\mathcal M}\qquad(5.1\text{-}11,\,12)
\]
En óptica casi siempre $\boldsymbol{\mathcal M}=0$ (medio no magnético). En el vacío ($\Pc=0$) cada componente del campo cumple la ecuación de onda $\nabla^2u-c_0^{-2}\partial_t^2u=0$, con $c_0=1/\sqrt{\epsilon_0\mu_0}$ (5.1-5, 6).

\textbf{Condiciones de frontera} (sin cargas ni corrientes libres en la interfaz): son continuas las componentes tangenciales $\mathcal E_t$, $\mathcal H_t$ y las normales $\mathcal D_n$, $\mathcal B_n$. En un conductor perfecto $\mathcal E_t=0$; por eso un espejo metálico refleja con desfase $\pi$. Estas condiciones son las que dan las fórmulas de Fresnel (§3.3).

\textbf{Energía: vector de Poynting.} El flujo de potencia por unidad de área es $\Sc=\Ec\times\Hc$ (5.1-13), y la intensidad es su promedio temporal, $I=|\langle\Sc\rangle|$. El teorema de Poynting es un balance de energía:
\[
\nabla\cdot\Sc=-\frac{\partial}{\partial t}\Big(\tfrac12\epsilon_0\mathcal E^2+\tfrac12\mu_0\mathcal H^2\Big)-\Ec\cdot\frac{\partial\Pc}{\partial t}-\mu_0\Hc\cdot\frac{\partial\boldsymbol{\mathcal M}}{\partial t}
\]
Se lee así: la potencia que sale de un volumen es igual a lo que pierde la energía almacenada en el campo, menos la potencia que el campo entrega a los dipolos del medio (los dos últimos términos). Esos términos son los responsables de la absorción.

La luz también lleva momento, con densidad $\epsilon_0\Ec\times\Bc=\Sc/c^2$. La presión de radiación es $\langle S\rangle/c$ sobre una superficie que absorbe y $2\langle S\rangle/c$ sobre una que refleja totalmente.

\subsection{Clasificación de los medios dieléctricos (§5.2)}
Un medio se describe por la relación entre $\Pc$ y $\Ec$. Cinco propiedades independientes:
\begin{itemize}
\item \textbf{Lineal}: $\Pc$ es proporcional a $\Ec$, así que vale la superposición.
\item \textbf{No dispersivo}: la respuesta es instantánea; $\Pc(t)$ depende solo de $\Ec(t)$.
\item \textbf{Homogéneo}: la relación no depende de la posición $\rv$.
\item \textbf{Isótropo}: la relación no depende de la dirección de $\Ec$, así que $\Pc\parallel\Ec$.
\item \textbf{Espacialmente no dispersivo}: $\Pc$ en $\rv$ depende solo de $\Ec$ en $\rv$ (respuesta local).
\end{itemize}

Lo que cambia según el caso:
\begin{itemize}
\item \textbf{Lineal, no dispersivo, homogéneo e isótropo}: $\Pc=\epsilon_0\chi\Ec$, con $\chi$ la susceptibilidad eléctrica. Entonces $\Dc=\epsilon\Ec$ con $\epsilon=\epsilon_0(1+\chi)$, y el índice es $n=\sqrt{\epsilon/\epsilon_0}=\sqrt{1+\chi}$ (5.2-1, 3, 13). La velocidad es $c=1/\sqrt{\epsilon\mu_0}=c_0/n$ y la densidad de energía $\mathcal W=\tfrac12\epsilon\mathcal E^2+\tfrac12\mu_0\mathcal H^2$.
\item \textbf{Inhomogéneo}, $\epsilon(\rv)$: $\nabla^2\Ec+\nabla\big(\tfrac1\epsilon\nabla\epsilon\cdot\Ec\big)-\mu_0\epsilon\,\partial_t^2\Ec=0$ (5.2-18). Si $\epsilon$ varía despacio frente a $\lambda$, el término central es despreciable y queda $\nabla^2\Ec-\partial_t^2\Ec/c^2(\rv)\approx0$. Una perturbación pequeña del índice, $\Delta n$, actúa como una fuente de radiación $\Delta\Pc=2\epsilon_0n\,\Delta n\,\Ec$ (5.2-20).
\item \textbf{Anisótropo}: cada componente de $\Pc$ depende de todas las de $\Ec$: $\mathcal P_i=\epsilon_0\sum_j\chi_{ij}\mathcal E_j$ (5.2-21). Es la base de los cristales (tema 3).
\item \textbf{Dispersivo}: el medio tiene memoria. $\Pc(t)=\epsilon_0\!\int x(t-t')\,\Ec(t')\,dt'$ (5.2-23), con $x(t)$ la respuesta al impulso. Por causalidad $x(t)=0$ para $t<0$. En frecuencia, $\chi(\nu)=\mathcal F\{x(t)\}$: cada frecuencia ve un índice distinto.
\item \textbf{No lineal}: $\Pc=\Psi(\Ec)$ con $\Psi$ no lineal, así que ya no vale la superposición. Es el objeto del curso (tema 4).
\end{itemize}

\begin{clave}[title={Ecuación de onda general (medio homogéneo e isótropo, lineal o no, dispersivo o no)}]
\[
\nabla^2\Ec-\frac{1}{c_0^2}\frac{\partial^2\Ec}{\partial t^2}=\mu_0\frac{\partial^2\Pc}{\partial t^2}\qquad(5.2\text{-}25)
\]
\textbf{Derivación}: tomar el rotacional de la ley de Faraday, sustituir la ley de Ampère, usar $\nabla\times\nabla\times\Ec=\nabla(\nabla\cdot\Ec)-\nabla^2\Ec$ y $\nabla\cdot\Ec=0$.

\textbf{Uso en óptica no lineal}: se separa $\Pc=\epsilon_0\chi\Ec+\Pc_{NL}$. La parte lineal pasa al lado izquierdo y cambia $c_0$ por $c=c_0/n$:
\[
\nabla^2\Ec-\frac{n^2}{c_0^2}\frac{\partial^2\Ec}{\partial t^2}=\mu_0\frac{\partial^2\Pc_{NL}}{\partial t^2}
\]
La polarización no lineal es una \emph{fuente} que radia campo a frecuencias nuevas.
\end{clave}

\subsection{Ondas monocromáticas (§5.3)}
Si todos los campos oscilan a una sola frecuencia, $\Ec=\mathrm{Re}\{\E(\rv)e^{j\omega t}\}$, cada derivada temporal se reemplaza por $j\omega$:
\[
\nabla\times\Hv=j\omega\D,\qquad \nabla\times\E=-j\omega\mathbf B .
\]
\textbf{Poynting complejo}: $\mathbf S=\tfrac12\E\times\Hv^*$, y la intensidad promedio es $\langle\Sc\rangle=\mathrm{Re}\{\mathbf S\}$ (5.3-10).

Cada componente cumple la \textbf{ecuación de Helmholtz} $\nabla^2U+k^2U=0$, con $k=nk_0=\omega\sqrt{\epsilon\mu_0}$ (5.3-16). En un medio dispersivo basta con dejar que todo dependa de la frecuencia: $\mathbf P=\epsilon_0\chi(\nu)\E$, $\epsilon(\nu)=\epsilon_0[1+\chi(\nu)]$, $k=\omega\sqrt{\epsilon(\nu)\mu_0}$.

\subsection{Ondas electromagnéticas elementales (§5.4)}
\textbf{Onda plana TEM.} $\E=\E_0e^{-j\kv\cdot\rv}$. Al sustituir en Maxwell salen $\kv\times\Hv_0=-\omega\epsilon\E_0$ y $\kv\times\E_0=\omega\mu_0\Hv_0$. Por tanto $\E$, $\Hv$ y $\kv$ son mutuamente perpendiculares (onda transversal electromagnética). La razón entre los campos es la \textbf{impedancia} del medio:
\[
\eta=\frac{E_0}{H_0}=\sqrt{\frac{\mu_0}{\epsilon}}=\frac{\eta_0}{n}\qquad(5.4\text{-}5,\,6)
\]

\begin{clave}[title={Intensidad y energía de una onda plana (5.4-8 a 5.4-10)}]
\[
I=\frac{|E_0|^2}{2\eta}=\frac{n|E_0|^2}{2\eta_0}=\tfrac12 n\epsilon_0c_0|E_0|^2,\qquad W=\tfrac12\epsilon|E_0|^2,\qquad I=cW
\]
Ejemplo: $10$ W/cm² en el vacío corresponde a $E_0\approx87$ V/cm. Campo magnético real: $\Hc=\frac1{\eta}\,\hat{\kv}\times\Ec$.
\end{clave}

\textbf{Dipolo oscilante} (lejos de él, $kr\gg1$): $\E\approx E_0\sin\theta\,U\,\hat{\boldsymbol\theta}$ y $\Hv\approx H_0\sin\theta\,U\,\hat{\boldsymbol\phi}$, con $U=e^{-jkr}/r$. El patrón de intensidad va como $\sin^2\theta$: no radia a lo largo de su eje. Es el modelo de cómo radia cada dipolo del medio inducido por $\Pc$.

\textbf{Haz gaussiano vectorial}: $\E=E_0\big(-\hat{\mathbf x}+\frac{x}{z+jz_0}\hat{\mathbf z}\big)U$ (5.4-20). Fuera del eje aparece una componente longitudinal ($z$): un haz de ancho finito no es exactamente transversal.

\textbf{Relación entre la teoría escalar y la electromagnética}: se identifica $U=E/\sqrt{2\eta}$ para que las intensidades coincidan. La teoría escalar falla cuando importa la polarización (Fresnel, guías de onda), porque $U$ se vuelve discontinuo donde cambia $\eta$.

\textbf{Haces vectoriales} (azimutal y radial): tienen forma de dona, con un nulo en el eje. Un haz con polarización radial enfocado da un foco más pequeño que uno con polarización lineal.

\subsection{Óptica ondulatoria escalar (§2.2)}
En la teoría escalar la luz se describe con una función de onda $u=\mathrm{Re}\{U(\rv)e^{j2\pi\nu t}\}$, con intensidad $I=|U|^2$ (2.2-10). $U$ cumple la ecuación de Helmholtz (2.2-7). Las soluciones básicas son:
\begin{itemize}
\item \textbf{Onda plana}: $U=Ae^{-j\kv\cdot\rv}$, frentes de onda planos separados $\lambda$.
\item \textbf{Onda esférica}: $U=\frac{A_0}{r}e^{-jkr}$ (2.2-15). La intensidad $I=|A_0|^2/r^2$ cae con el cuadrado de la distancia.
\item \textbf{Aproximación de Fresnel} (onda paraboloidal): cerca del eje ($\theta\ll1$), $r\approx z+\frac{x^2+y^2}{2z}$, y
\[
U\approx\frac{A_0}{z}e^{-jkz}\exp\!\Big[-jk\frac{x^2+y^2}{2z}\Big]\qquad(2.2\text{-}17)
\]
Vale si $N_F\theta_m^2/4\ll1$, donde $N_F=a^2/\lambda z$ es el número de Fresnel ($a$ el radio de la zona observada).
\item \textbf{Onda paraxial}: una onda plana modulada por una envolvente que varía despacio, $U=A(\rv)e^{-jkz}$, con $|\partial_zA|\ll k|A|$ y $|\partial_z^2A|\ll k^2|A|$ (la envolvente apenas cambia en una longitud de onda).
\end{itemize}
\begin{clave}[title={Ecuación de Helmholtz paraxial}]
\[
\nabla_T^2A-j2k\,\frac{\partial A}{\partial z}=0\qquad(2.2\text{-}23),\qquad \nabla_T^2=\partial_x^2+\partial_y^2
\]
Se obtiene al despreciar $\partial_z^2A$. Sus soluciones son la onda paraboloidal, el haz gaussiano (tema 2) y los haces de Hermite–Gauss.
\end{clave}

%=====================================================================
\section{Haz gaussiano; absorción y dispersión; pulsos}
%=====================================================================

\begin{idea}
Tres ideas. (1) El haz gaussiano es la forma real de un haz láser: se enfoca hasta un radio mínimo y luego diverge, y la intensidad en el foco es la que controla los efectos no lineales. (2) La susceptibilidad es compleja: la parte real da el índice y la imaginaria da la absorción, y las dos están ligadas por la causalidad (Kramers–Kronig). (3) Como el índice depende de la frecuencia, cada componente espectral de un pulso viaja a distinta velocidad y el pulso se ensancha.
\end{idea}

\subsection{Haz gaussiano (Saleh §3.1)}
Se obtiene de la onda paraboloidal reemplazando $z$ por el parámetro complejo $q(z)=z+jz_0$:
\[
A=\frac{A_1}{q}\exp\!\Big[-jk\frac{\rho^2}{2q}\Big],\qquad \frac1{q(z)}=\frac1{R(z)}-j\frac{\lambda}{\pi W^2(z)}\qquad(3.1\text{-}6)
\]
La parte real de $1/q$ da la curvatura del frente de onda y la imaginaria da el ancho del haz.

\begin{clave}[title={Amplitud compleja del haz gaussiano (3.1-7)}]
\[
U(\rho,z)=A_0\frac{W_0}{W(z)}\exp\!\Big[-\frac{\rho^2}{W^2(z)}\Big]\exp\!\Big[-jkz-jk\frac{\rho^2}{2R(z)}+j\zeta(z)\Big]
\]
\[
W(z)=W_0\sqrt{1+(z/z_0)^2},\qquad R(z)=z\big[1+(z_0/z)^2\big],\qquad \zeta(z)=\tan^{-1}(z/z_0),
\]
\[
W_0=\sqrt{\lambda z_0/\pi},\qquad z_0=\pi W_0^2/\lambda\qquad(\lambda=\lambda_0/n).
\]
$W(z)$: radio del haz (donde la amplitud cae a $1/e$). $W_0$: cintura (radio mínimo, en $z=0$). $z_0$: rango de Rayleigh. $R(z)$: radio de curvatura del frente de onda. $\zeta(z)$: fase de Gouy.
\end{clave}

Propiedades:
\begin{itemize}
\item \textbf{Intensidad}: $I(\rho,z)=I_0\big(\frac{W_0}{W}\big)^2e^{-2\rho^2/W^2}=\frac{2P}{\pi W^2}e^{-2\rho^2/W^2}$, con potencia total $P=\frac12I_0\pi W_0^2$. En el eje, $I=I_0/[1+(z/z_0)^2]$.
\item \textbf{Potencia dentro de un círculo} de radio $\rho_0$: fracción $1-e^{-2\rho_0^2/W^2}$ (86\,\% dentro de $W$; 99\,\% dentro de $1.5W$).
\item \textbf{Divergencia}: lejos del foco el haz se abre con semiángulo $\theta_0=\lambda/\pi W_0$ (3.1-20). Cuanto más pequeña la cintura, más diverge. La \textbf{profundidad de foco} es $2z_0=2\pi W_0^2/\lambda$.
\item \textbf{En $z=z_0$}: $W=\sqrt2W_0$ (el área se duplica) y la intensidad en el eje es la mitad del pico. Ahí la curvatura es máxima, $R_{\min}=2z_0$. En la cintura el frente es plano, $R(0)=\infty$.
\item \textbf{Fase de Gouy}: la fase en el eje es $kz-\zeta(z)$. El término $\zeta$ es un retraso extra respecto a una onda plana, con total $\pi$ al pasar de $z=-\infty$ a $+\infty$; vale $\zeta(z_0)=\pi/4$.
\item \textbf{Problema inverso}: si en un punto se conocen $W$ y $R$, la cintura y su posición son
$z=\dfrac{R}{1+(\lambda R/\pi W^2)^2}$, \ $W_0=\dfrac{W}{\sqrt{1+(\pi W^2/\lambda R)^2}}$ (3.1-25, 26).
\item \textbf{Calidad del haz}: $M^2=\theta_m/\theta_0=\dfrac{2W_m\cdot2\theta_m}{4\lambda/\pi}\ge1$. $M^2=1$ es un gaussiano ideal; un haz real diverge $M^2$ veces más.
\item \textbf{Lentes y sistemas ópticos} (§3.2): una lente delgada transforma $\frac1{q_2}=\frac1{q_1}-\frac1f$. Un sistema con matriz ABCD: $q_2=\frac{Aq_1+B}{Cq_1+D}$. Un haz colimado ($z_0\gg f$) se enfoca a una cintura $W_0'\approx\frac{\lambda}{\pi W_0}f=\theta_0f$.
\end{itemize}

\subsection{Absorción (Saleh §5.5A)}
Cuando el medio absorbe, la susceptibilidad es compleja, $\chi=\chi'+j\chi''$, y también lo es el número de onda:
\[
k=k_0\sqrt{1+\chi}=\beta-j\frac{\alpha}{2}.
\]
La onda queda $U\propto e^{-\alpha z/2}e^{-j\beta z}$: la parte real $\beta=nk_0$ da la fase (el índice) y la parte imaginaria la atenuación. La amplitud decae con $\alpha/2$ y la intensidad con $\alpha$, $I\propto e^{-\alpha z}$ ($\alpha$ es el coeficiente de absorción).

\begin{clave}[title={Índice y absorción a partir de $\chi$}]
\[
n-j\frac{\alpha}{2k_0}=\sqrt{1+\chi'+j\chi''}\qquad(5.5\text{-}5)\qquad(\text{raíz del 4.º cuadrante});\qquad \eta=\frac{\eta_0}{\sqrt{1+\chi}}
\]
\begin{itemize}[leftmargin=1.2em]
\item \textbf{Absorción débil} ($|\chi''|\ll1+\chi'$): $n\approx\sqrt{1+\chi'}$, \ $\alpha\approx-\dfrac{k_0}{n}\chi''$.
\item \textbf{Absorción fuerte} ($|\chi''|\gg|1+\chi'|$): $n\approx\sqrt{-\chi''/2}$, \ $\alpha\approx2k_0\sqrt{-\chi''/2}$.
\item \textbf{Absorbente diluido en un huésped} de índice $n_0$: $\chi_{\rm tot}=n_0^2-1+\chi$, y entonces $n\approx n_0+\dfrac{\chi'}{2n_0}$, \ $\alpha\approx-\dfrac{k_0\chi''}{n_0}$.
\end{itemize}
\end{clave}
En la convención del Saleh, $\chi''<0$ significa absorción. Si $\chi''>0$ el medio \textbf{amplifica}: $\alpha<0$ y se define la ganancia $\gamma=-\alpha$ (medio láser). Atenuación en dB: dB/km $=4.343\,\alpha$ [km$^{-1}$].

\subsection{Relaciones de Kramers–Kronig (§5.5B)}
Si el sistema es lineal, causal, invariante en el tiempo y de respuesta real ($\chi(-\nu)=\chi^*(\nu)$), las partes real e imaginaria de $\chi$ no son independientes: conocida una en todo el espectro, se calcula la otra.
\[
\chi'(\nu)=\frac2\pi\int_0^\infty\frac{s\,\chi''(s)}{s^2-\nu^2}\,ds,\qquad
\chi''(\nu)=\frac2\pi\int_0^\infty\frac{\nu\,\chi'(s)}{\nu^2-s^2}\,ds\qquad(\text{valor principal})
\]
\textbf{Consecuencia}: todo medio dispersivo ($\chi'$ varía con $\nu$) es también absorbente en alguna región del espectro.

La dispersión de los vidrios se caracteriza con el \textbf{número de Abbe} $\mathbb V=(n_d-1)/(n_F-n_C)$, con las líneas $d$, $F$, $C$ en $587.6$, $486.1$ y $656.3$ nm. Un $\mathbb V$ grande significa poca dispersión.

\subsection{Medio resonante: modelo de Lorentz (§5.5C)}
Cada electrón ligado se modela como un oscilador amortiguado forzado por el campo: $\ddot x+\sigma\dot x+\omega_0^2x=F/m$, con $F=-e\mathcal E$. Con $N$ electrones por unidad de volumen, $\mathcal P=-Nex$, y la polarización cumple
\[
\ddot{\mathcal P}+\sigma\dot{\mathcal P}+\omega_0^2\mathcal P=\omega_0^2\epsilon_0\chi_0\mathcal E,\qquad \chi_0=\frac{Ne^2}{\epsilon_0m\omega_0^2}\quad(5.5\text{-}17)
\]
($\chi_0$ es la susceptibilidad a frecuencia cero.) Para una onda armónica sale:

\begin{clave}[title={Susceptibilidad de Lorentz (5.5-19)}]
\[
\chi(\nu)=\chi_0\frac{\nu_0^2}{\nu_0^2-\nu^2+j\nu\Delta\nu},\qquad \Delta\nu=\frac{\sigma}{2\pi}
\]
\[
\chi'=\chi_0\frac{\nu_0^2(\nu_0^2-\nu^2)}{(\nu_0^2-\nu^2)^2+(\nu\Delta\nu)^2},\qquad
\chi''=-\chi_0\frac{\nu_0^2\,\nu\Delta\nu}{(\nu_0^2-\nu^2)^2+(\nu\Delta\nu)^2}
\]
\textbf{Cerca de la resonancia} ($\nu\approx\nu_0$): $\chi\approx\chi_0\dfrac{\nu_0/2}{(\nu_0-\nu)+j\Delta\nu/2}$, de modo que
\[
\chi''\approx-\chi_0\frac{\nu_0\Delta\nu}{4}\,\frac{1}{(\nu_0-\nu)^2+(\Delta\nu/2)^2}\quad(\text{lorentziana de ancho FWHM }\Delta\nu),\qquad \chi'=2\frac{\nu-\nu_0}{\Delta\nu}\chi''.
\]
\end{clave}
Interpretación: la absorción ($-\chi''$) es un pico en $\nu_0$ con altura $-\chi''(\nu_0)=\chi_0Q$, donde $Q=\nu_0/\Delta\nu$. El índice ($\chi'$) cambia de signo en la resonancia, con extremos $\approx\pm\chi_0Q/2$ en $\nu_0(1\mp\frac1{2Q})$. Lejos de la resonancia $\chi\approx\chi_0\nu_0^2/(\nu_0^2-\nu^2)$ es real: el medio es transparente y solo dispersa.

Diluido en un huésped: $\alpha(\nu)\approx-\dfrac{2\pi\nu}{n_0c_0}\chi''$, \ $n\approx n_0+\dfrac{\chi'}{2n_0}$.

\textbf{Ecuación de Sellmeier}: un material real tiene varias resonancias; lejos de todas ellas,
\[
n^2=1+\sum_i\chi_{0i}\frac{\lambda^2}{\lambda^2-\lambda_i^2}\qquad(5.5\text{-}28).
\]
Sílice: $\lambda_i=0.06840,\ 0.1162,\ 9.8962$ µm; \ $\chi_{0i}=0.6962,\ 0.4079,\ 0.8975$.

\subsection{Medios conductores (§5.5D)}
Una conductividad $\sigma$ entra como una permitividad efectiva compleja, $\epsilon_e=\epsilon+\sigma/j\omega$. En un buen conductor ($\sigma\gg\omega\epsilon$):
\[
n\approx\sqrt{\sigma/2\omega\epsilon_0},\qquad \alpha\approx\sqrt{2\omega\mu_0\sigma},\qquad d_p=1/\alpha\ (\text{profundidad de penetración}).
\]
Ejemplo: cobre a 1 µm, $d_p=1.9$ nm; la luz casi no penetra en un metal.

\textbf{Modelo de Drude} (electrones libres con tiempo de colisión $\tau$): $\sigma=\dfrac{\sigma_0}{1+j\omega\tau}$, con $\sigma_0=Ne^2\tau/m$. A alta frecuencia, $\epsilon_e=\epsilon_0\big(1-\omega_p^2/\omega^2\big)$, con frecuencia de plasma $\omega_p=\sqrt{Ne^2/\epsilon_0m}$.
\begin{itemize}
\item $\omega<\omega_p$: $\epsilon_e<0$, la onda es evanescente y el metal refleja.
\item $\omega>\omega_p$: el metal se vuelve transparente, con $n=\sqrt{1-\omega_p^2/\omega^2}<1$.
\end{itemize}

\subsection{Pulsos en medios dispersivos (§5.6; §22.3)}
Un pulso contiene una banda de frecuencias alrededor de $\omega_0$. Se expande la constante de propagación $\beta(\omega)=\omega n(\omega)/c_0$ en serie de Taylor, con $\Omega=\omega-\omega_0$:
\[
\beta(\omega)\approx\beta_0+\beta'\Omega+\tfrac12\beta''\Omega^2 .
\]
Cada término tiene un papel: $\beta_0$ da la fase de la portadora, $\beta'$ el retardo de la envolvente (velocidad de grupo) y $\beta''$ el ensanchamiento (dispersión de la velocidad de grupo).

\begin{clave}[title={Velocidad de grupo y coeficientes de dispersión}]
Velocidad de grupo $v=1/\beta'$; índice de grupo $N=n-\lambda_0\dfrac{dn}{d\lambda_0}=n+\omega\dfrac{dn}{d\omega}$; \ $v=c_0/N$; retardo $\tau_d=z/v=Nz/c_0$.
\[
D_\nu=\frac{d}{d\nu}\Big(\frac1v\Big)=2\pi\beta''=\frac{\lambda_0^3}{c_0^2}\frac{d^2n}{d\lambda_0^2}\ \ [\mathrm{s/(m\cdot Hz)}],\qquad
D_\lambda=-\frac{\lambda_0}{c_0}\frac{d^2n}{d\lambda_0^2}=-\frac{c_0}{\lambda_0^2}D_\nu\ \ [\mathrm{ps/(km\cdot nm)}]
\]
Ensanchamiento de un pulso emitido por una fuente de ancho espectral grande: $\sigma_\tau=|D_\nu|\sigma_\nu z=|D_\lambda|\sigma_\lambda z$.
\end{clave}
\begin{itemize}
\item \textbf{Dispersión normal}: $\beta''>0$, $D_\nu>0$, $D_\lambda<0$. El azul viaja más lento que el rojo y llega después.
\item \textbf{Dispersión anómala}: $\beta''<0$, $D_\lambda>0$; el azul llega antes.
\item Sílice: $D_\lambda=0$ en $1.276$ µm (mínimo de $N$); fibra estándar $\approx1.312$ µm; $D_\lambda\approx+17$ ps/(km·nm) a 1.55 µm (anómala).
\item Notación de Agrawal (fibras): $\beta_2=\beta''$, \ $D_\lambda=-\dfrac{2\pi c_0}{\lambda_0^2}\beta_2$.
\end{itemize}

\textbf{Ecuación de la envolvente.} En un marco que viaja con el pulso ($t\to t-z/v$), la envolvente $A(z,t)$ cumple
\[
\frac{\partial A}{\partial z}=j\frac{\beta''}{2}\frac{\partial^2A}{\partial t^2}=j\frac{D_\nu}{4\pi}\frac{\partial^2A}{\partial t^2}.
\]
Es la misma forma que la ecuación paraxial: la dispersión ensancha el pulso en el tiempo igual que la difracción ensancha el haz en el espacio. En frecuencia la solución es $A(z,\nu)=A(0,\nu)e^{-j\pi D_\nu\nu^2z}$: un filtro de fase pura que no cambia el espectro $|A(\nu)|^2$, solo las fases relativas.

\begin{clave}[title={Pulso gaussiano en un medio dispersivo}]
Entrada $A(0,t)=e^{-t^2/\tau_0^2}$, con $\tau_0$ el semiancho a $1/e$ de la amplitud ($\tau_{\rm FWHM}=\sqrt{2\ln2}\,\tau_0=1.177\tau_0$).
\[
\tau(z)=\tau_0\sqrt{1+(z/z_0)^2},\qquad z_0=\frac{\pi\tau_0^2}{|D_\nu|}=\frac{\tau_0^2}{2|\beta''|};\qquad \text{chirp }a(z)=\pm z/z_0;\qquad \tau\approx\frac{|D_\nu|z}{\pi\tau_0}\ (z\gg z_0).
\]
Notación de Agrawal: $e^{-T^2/2T_0^2}$ y longitud de dispersión $L_D=T_0^2/|\beta_2|$ (con $\tau_0=\sqrt2T_0$ es el mismo $z_0$). Un gaussiano sin chirp cumple $\tau_{\rm FWHM}\Delta\nu=0.44$ (límite de transformada).
\end{clave}
El pulso adquiere \emph{chirp}: su frecuencia instantánea cambia a lo largo del pulso. Cuanto más corto es el pulso de entrada, más rápido se ensancha ($z_0\propto\tau_0^2$).

\textbf{Luz rápida y lenta}: cerca de una resonancia, $N$ puede ser $\gg1$, menor que 1 o incluso negativo. Aun así, la velocidad a la que viaja la información nunca supera $c_0$.

%=====================================================================
\section{Polarización y cristales (Saleh cap.~6)}
%=====================================================================

\begin{idea}
La polarización es la dirección en que oscila $\Ec$. En los cristales anisótropos el índice depende de esa dirección, así que dos polarizaciones viajan a velocidades distintas (birrefringencia). Esto importa para el curso porque la birrefringencia es la herramienta estándar para lograr el ajuste de fase en la generación de segundo armónico, y porque $d_{\rm eff}$ depende de las polarizaciones de las ondas.
\end{idea}

\subsection{Estados de polarización (§6.1A)}
Una onda plana hacia $+z$: $\Ec=\mathrm{Re}\{(A_x\hat x+A_y\hat y)e^{j\omega(t-z/c)}\}$, con $A_{x,y}=\mathsf a_{x,y}e^{j\varphi_{x,y}}$. Lo que importa es el desfase $\varphi=\varphi_y-\varphi_x$ y la razón de amplitudes $\mathsf r=\mathsf a_y/\mathsf a_x$. En general la punta de $\Ec$ dibuja una elipse:
\[
\frac{\mathcal E_x^2}{\mathsf a_x^2}+\frac{\mathcal E_y^2}{\mathsf a_y^2}-2\cos\varphi\,\frac{\mathcal E_x\mathcal E_y}{\mathsf a_x\mathsf a_y}=\sin^2\varphi\qquad(6.1\text{-}5)
\]
La elipse se caracteriza por su orientación $\psi$ y su elipticidad $\chi$ ($\tan\chi=b/a$, razón de semiejes):
\[
\tan2\psi=\frac{2\mathsf r}{1-\mathsf r^2}\cos\varphi,\qquad \sin2\chi=\frac{2\mathsf r}{1+\mathsf r^2}\sin\varphi .
\]
Casos particulares: \textbf{lineal} si $\varphi=0$ o $\pi$; \textbf{circular} si $\mathsf a_x=\mathsf a_y$ y $\varphi=\pm\pi/2$. Con $\varphi=+\pi/2$ es circular \textbf{derecha} (gira en sentido horario visto de frente, mirando hacia la fuente); con $-\pi/2$, izquierda.

\textbf{Parámetros de Stokes}: describen la polarización con cantidades medibles (intensidades):
\[
S_0=|A_x|^2+|A_y|^2,\quad S_1=|A_x|^2-|A_y|^2,\quad S_2=2\mathrm{Re}\{A_x^*A_y\}=|A_{45}|^2-|A_{135}|^2,\quad S_3=2\mathrm{Im}\{A_x^*A_y\}=|A_R|^2-|A_L|^2,
\]
con $S_1^2+S_2^2+S_3^2=S_0^2$ para luz totalmente polarizada. $S_1$ compara horizontal con vertical, $S_2$ las diagonales y $S_3$ las circulares.

\textbf{Esfera de Poincaré}: el punto $(S_1,S_2,S_3)/S_0=(\cos2\chi\cos2\psi,\ \cos2\chi\sin2\psi,\ \sin2\chi)$ representa cada estado. El ecuador contiene los estados lineales y los polos los circulares (norte: derecha).

\subsection{Matrices de Jones (§6.1B)}
La polarización se representa con el vector de Jones $\J=(A_x,A_y)^T$; la intensidad es $I=(|A_x|^2+|A_y|^2)/2\eta$. Cada elemento óptico es una matriz $2\times2$ que actúa sobre $\J$.

Vectores básicos: lineal a ángulo $\theta$: $(\cos\theta,\sin\theta)$; circular derecha: $\frac1{\sqrt2}(1,j)$; circular izquierda: $\frac1{\sqrt2}(1,-j)$. Dos estados son ortogonales si $A_{1x}A_{2x}^*+A_{1y}A_{2y}^*=0$.

Cambio de base: $A_R=\frac{A_x-jA_y}{\sqrt2}$, $A_L=\frac{A_x+jA_y}{\sqrt2}$; \ $A_{45}=\frac{A_x+A_y}{\sqrt2}$, $A_{135}=\frac{A_y-A_x}{\sqrt2}$.

\begin{clave}[title={Matrices de Jones de los elementos básicos}]
\renewcommand{\arraystretch}{1.6}
\begin{tabular}{@{}ll@{}}
Polarizador lineal en $x$ & $\begin{pmatrix}1&0\\0&0\end{pmatrix}$\\
Polarizador lineal a ángulo $\theta$ & $\begin{pmatrix}\cos^2\theta&\sin\theta\cos\theta\\ \sin\theta\cos\theta&\sin^2\theta\end{pmatrix}$\\
Retardador de onda (eje rápido en $x$) & $\begin{pmatrix}1&0\\0&e^{-j\Gamma}\end{pmatrix}$\\
Rotador de polarización & $\begin{pmatrix}\cos\theta&-\sin\theta\\ \sin\theta&\cos\theta\end{pmatrix}$
\end{tabular}

\textbf{Rotación de ejes}: $\J'=\R(\theta)\J$ con $\R(\theta)=\begin{psmallmatrix}\cos\theta&\sin\theta\\-\sin\theta&\cos\theta\end{psmallmatrix}$. Para un elemento girado: $\T'=\R(\theta)\T\R(-\theta)$ y $\T=\R(-\theta)\T'\R(\theta)$.

\textbf{Cascada}: $\T=\T_N\cdots\T_2\T_1$ (el primer elemento que atraviesa la luz va a la derecha).

\textbf{Modos normales}: los autovectores de $\T$ ($\T\J=\mu\J$) son los estados que atraviesan el elemento sin cambiar de polarización.
\end{clave}
Resultados útiles:
\begin{itemize}
\item Lámina $\lambda/4$ ($\Gamma=\pi/2$): convierte luz lineal a 45° de los ejes en circular, y viceversa.
\item Lámina $\lambda/2$ ($\Gamma=\pi$): convierte lineal a $\theta$ en lineal a $-\theta$ respecto al eje rápido (gira la polarización $2\theta$), e intercambia circular derecha e izquierda.
\item Ley de Malus: $I=I_0\cos^2\theta$. Con $N$ polarizadores a pasos de $\pi/2N$, la amplitud transmitida es $\cos^N(\pi/2N)\to1$ cuando $N\to\infty$: se puede girar la polarización 90° casi sin pérdidas.
\item Retardador entre polarizadores cruzados, con sus ejes a 45°: transmitancia $\mathcal T=\sin^2(\Gamma/2)$ (6.6-1). Es el principio del modulador electroóptico.
\end{itemize}

\subsection{Reflexión y refracción (§6.2)}
En una interfaz plana entre índices $n_1$ y $n_2$ vale la ley de Snell, $n_1\sin\theta_1=n_2\sin\theta_2$. La reflexión depende de la polarización, que se separa en dos casos: \textbf{TE} (o s), con $\E$ perpendicular al plano de incidencia, y \textbf{TM} (o p), con $\E$ en el plano de incidencia.

\begin{clave}[title={Ecuaciones de Fresnel (convención del Saleh)}]
\[
r_{TE}=\frac{n_1\cos\theta_1-n_2\cos\theta_2}{n_1\cos\theta_1+n_2\cos\theta_2},\qquad t_{TE}=1+r_{TE}
\]
\[
r_{TM}=\frac{n_1\sec\theta_1-n_2\sec\theta_2}{n_1\sec\theta_1+n_2\sec\theta_2},\qquad t_{TM}=(1+r_{TM})\frac{\cos\theta_1}{\cos\theta_2}
\]
Reflectancia $\mathcal R=|r|^2$; transmitancia $\mathcal T=\dfrac{n_2\cos\theta_2}{n_1\cos\theta_1}|t|^2=1-\mathcal R$. A incidencia normal: $\mathcal R=\Big(\dfrac{n_1-n_2}{n_1+n_2}\Big)^2$.
\end{clave}
\begin{itemize}
\item \textbf{Ángulo de Brewster} $\theta_B=\tan^{-1}(n_2/n_1)$: la onda TM no se refleja ($r_{TM}=0$) y $\theta_1+\theta_2=90^\circ$. La luz reflejada a ese ángulo sale polarizada TE.
\item \textbf{Reflexión total interna} (de más a menos denso, $n_1>n_2$): por encima del ángulo crítico $\theta_c=\sin^{-1}(n_2/n_1)$, $|r|=1$ y solo cambia la fase:
\[
\tan\frac{\varphi_{TE}}2=\sqrt{\frac{\cos^2\theta_c}{\cos^2\theta_1}-1},\qquad \tan\frac{\varphi_{TM}}2=-\frac{1}{\sin^2\theta_c}\sqrt{\frac{\cos^2\theta_c}{\cos^2\theta_1}-1}.
\]
Al otro lado queda una onda evanescente que decae como $e^{-\gamma y}$, con $\gamma=k_0\sqrt{n_1^2\sin^2\theta_1-n_2^2}$. Es el mecanismo de guiado en guías de onda y fibras.
\item Valores: vidrio refleja 4\,\%, GaAs 32\,\%. Una lámina con dos caras: $\mathcal R(1+\mathcal T^2)$. Un conductor (Hagen–Rubens): $\mathcal R\approx1-2\sqrt{2\epsilon_0\omega/\sigma}$.
\end{itemize}

\subsection{Medios anisótropos (§6.3)}
En un cristal anisótropo $\D=\boldsymbol\epsilon\E$, con $\boldsymbol\epsilon$ un tensor simétrico. Siempre existen tres \textbf{ejes principales} en los que el tensor es diagonal: $D_i=\epsilon_iE_i$, y cada eje tiene su índice $n_i=\sqrt{\epsilon_i/\epsilon_0}$. Según cuántos índices distintos hay:
\begin{itemize}
\item \textbf{Biaxial}: tres índices distintos.
\item \textbf{Uniaxial}: $n_1=n_2=n_o$ (ordinario) y $n_3=n_e$ (extraordinario). El eje $z$ es el \textbf{eje óptico}. Es positivo si $n_e>n_o$ y negativo si $n_e<n_o$.
\item \textbf{Isótropo}: un solo índice (cristales cúbicos).
\end{itemize}

\begin{clave}[title={Elipsoide de índices y cristal uniaxial}]
\textbf{Elipsoide de índices}: $\dfrac{x_1^2}{n_1^2}+\dfrac{x_2^2}{n_2^2}+\dfrac{x_3^2}{n_3^2}=1$. \textbf{Receta}: para una onda que viaja en la dirección $\ud$, se corta el elipsoide con el plano normal a $\ud$. La elipse resultante tiene semiejes $n_a$ y $n_b$, que son los índices de los dos modos que pueden propagarse, y sus ejes dan las direcciones de $\D_a$ y $\D_b$.

\textbf{Uniaxial}, con $\ud$ a un ángulo $\theta$ del eje óptico:
\[
\frac{1}{n^2(\theta)}=\frac{\cos^2\theta}{n_o^2}+\frac{\sin^2\theta}{n_e^2}\qquad(6.3\text{-}15)
\]
\begin{itemize}[leftmargin=1.2em]
\item \textbf{Onda ordinaria (o)}: índice $n_o$ para cualquier $\theta$; $\D$ perpendicular al plano que forman $\kv$ y el eje óptico; $\E\parallel\D$ y la energía viaja paralela a $\kv$.
\item \textbf{Onda extraordinaria (e)}: índice $n(\theta)$, entre $n_o$ ($\theta=0$) y $n_e$ ($\theta=90^\circ$); $\D$ en ese plano; $\E\not\parallel\D$, de modo que la energía ($\Sc$) no viaja paralela a $\kv$.
\end{itemize}
\end{clave}

Geometría general: $\D\perp\kv$ y $\E\perp\Sc$; los cuatro vectores $\D,\E,\kv,\Sc$ son coplanares. La \textbf{superficie de números de onda} es $\sum_j\frac{u_j^2k^2}{k^2-n_j^2k_0^2}=1$; para el uniaxial se factoriza en una esfera (onda o) y un elipsoide (onda e):
\[
(k^2-n_o^2k_0^2)\Big(\frac{k_1^2+k_2^2}{n_e^2}+\frac{k_3^2}{n_o^2}-k_0^2\Big)=0 .
\]
La dirección del rayo (energía) es normal a esta superficie.

\textbf{Walk-off} (separación espacial entre $\kv$ y $\Sc$ en la onda e): con $\theta_S$ el ángulo de $\Sc$ respecto al eje,
\[
\tan\theta_S=\frac{n_o^2}{n_e^2}\tan\theta,\qquad \rho=\theta-\theta_S,\qquad
\tan\rho=\frac{(n_e^2-n_o^2)\sin\theta\cos\theta}{n_e^2\cos^2\theta+n_o^2\sin^2\theta}.
\]
El máximo es $\tan\rho_{\max}=\frac{n_e^2-n_o^2}{2n_en_o}$, en $\tan\theta=n_e/n_o$. En óptica no lineal el walk-off separa los haces fundamental y armónico y limita la longitud útil del cristal.

\textbf{Doble refracción}: en la superficie del cristal, la onda o cumple $\sin\theta_1=n_o\sin\theta_o$ y la onda e cumple $\sin\theta_1=n(\theta_a+\theta_e)\sin\theta_e$ ($\theta_a$: ángulo del eje óptico con la normal). Incluso a incidencia normal, con el eje oblicuo, el rayo e se desvía: así funcionan los divisores de polarización.

\textbf{Lámina retardadora}: una lámina de espesor $d$ con los dos índices $n_1$, $n_2$ introduce un retardo $\Gamma=(n_2-n_1)k_0d=2\pi\,\Delta n\,d/\lambda_0$. Ejemplo: mica a 633 nm, $d=63.3$ µm da $\lambda/2$.

\subsection{Actividad óptica, efecto Faraday y cristales líquidos (§6.4–6.6)}
\textbf{Actividad óptica}: en algunos materiales (cuarzo, soluciones de azúcar) los modos normales son las dos polarizaciones circulares, con índices $n_\pm=\sqrt{n^2\pm G}$. Constitutivamente, $\D=\boldsymbol\epsilon\E+j\epsilon_0\mathbf G\times\E$, con $\mathbf G=\xi\kv$. Como las dos circulares se desfasan, una polarización lineal gira a razón de
\[
\rho=\frac{\pi}{\lambda_0}(n_--n_+)\approx-\frac{\pi G}{\lambda_0n}\quad[\text{rad/m}].
\]
Es un efecto \textbf{recíproco}: al ir y volver, la rotación total es cero. Cuarzo: $\approx31^\circ$/mm a 500 nm.

\textbf{Efecto Faraday}: un campo magnético $\mathbf B$ a lo largo de la propagación produce la misma rotación, con $\rho=\mathfrak VB$ ($\mathfrak V$: constante de Verdet) y $\mathbf G=\gamma\mathbf B$. La diferencia es que es \textbf{no recíproco}: al ir y volver la rotación se duplica. Con esto se construye el \textbf{aislador óptico}: polarizador, rotador de Faraday de 45° y polarizador a 45°; la luz reflejada vuelve girada 90° y queda bloqueada.

\textbf{Cristal líquido nemático torcido}: las moléculas giran con la profundidad, $\theta=\alpha z$. Si la birrefringencia domina sobre la torsión ($\beta=(n_e-n_o)k_0\gg\alpha$), la polarización lineal sigue la torsión y el cristal actúa como rotador. Matriz: $\T=\R(-\alpha d)\,\mathrm{diag}(e^{-j\beta d/2},e^{j\beta d/2})$.

\textbf{Polarizadores}: dicroicos (Polaroid, rejilla de alambres), por ángulo de Brewster, y prismas birrefringentes (Wollaston, Rochon, Glan–Thompson).

%=====================================================================
\section{Susceptibilidad óptica no lineal (Boyd cap.~1)}
%=====================================================================

\begin{idea}
Con campos débiles los electrones oscilan en la zona parabólica del potencial y $\Pc$ es proporcional a $\Ec$. Con campos intensos (láseres) el potencial deja de ser parabólico y $\Pc$ adquiere términos en $\Ec^2$, $\Ec^3$, etc. Como $\cos^2\omega t$ contiene $\cos2\omega t$, el medio radia a frecuencias que no estaban en el haz incidente: segundo armónico, suma y diferencia de frecuencias, tercer armónico. El término cúbico también modifica el índice a la misma frecuencia (efecto Kerr). Las simetrías del material deciden qué términos sobreviven: sin $\chi^{(2)}$ en medios centrosimétricos, $\chi^{(3)}$ en todos.
\end{idea}

\textbf{Convención del Boyd}: $\tilde E(t)=Ee^{-i\omega t}+\cc=\sum_nE(\omega_n)e^{-i\omega_nt}$, con $E(-\omega)=E(\omega)^*$. Si el campo real es $\tilde E=\mathcal E\cos\omega t$, entonces $E=\mathcal E/2$. Intensidad: $I=2n\epsilon_0c|E|^2=\frac12n\epsilon_0c\mathcal E^2$.

\subsection{Expansión de la polarización y órdenes de magnitud (§1.1)}
\begin{clave}[title={Polarización no lineal (1.1.2)}]
\[
\tilde P(t)=\epsilon_0\big[\chi^{(1)}\tilde E(t)+\chi^{(2)}\tilde E^2(t)+\chi^{(3)}\tilde E^3(t)+\cdots\big]\equiv\tilde P^{(1)}+\tilde P^{(2)}+\tilde P^{(3)}+\cdots
\]
Esta forma escalar supone respuesta instantánea, lo que implica un medio sin pérdidas ni dispersión. Unidades SI: $\chi^{(1)}$ adimensional, $\chi^{(2)}$ en m/V, $\chi^{(3)}$ en m²/V².
\end{clave}

\textbf{¿Cuándo importan los términos no lineales?} La escala natural es el campo que liga al electrón con el núcleo en el átomo de hidrógeno:
\[
E_{\rm at}=\frac{e}{4\pi\epsilon_0a_0^2}=5.14\times10^{11}\ \mathrm{V/m}.
\]
Cada orden es menor que el anterior en un factor $\mathcal E/E_{\rm at}$: $P^{(n+1)}/P^{(n)}\sim\mathcal E/E_{\rm at}$. De ahí las estimaciones (fuera de resonancia)
\[
\chi^{(2)}\sim\frac{1}{E_{\rm at}}\approx1.94\times10^{-12}\ \mathrm{m/V},\qquad \chi^{(3)}\sim\frac{1}{E_{\rm at}^2}\approx3.78\times10^{-24}\ \mathrm{m^2/V^2}.
\]
La intensidad correspondiente es $I_{\rm at}=\frac12\epsilon_0cE_{\rm at}^2=3.5\times10^{20}$ W/m² $=3.5\times10^{16}$ W/cm². Por eso la óptica no lineal necesita láseres.

\textbf{Ecuación de onda} con fuente no lineal (1.1.5):
\[
\nabla^2\tilde E-\frac{n^2}{c^2}\frac{\partial^2\tilde E}{\partial t^2}=\frac{1}{\epsilon_0c^2}\frac{\partial^2\tilde P^{NL}}{\partial t^2}\qquad\Big(\frac1{\epsilon_0c^2}=\mu_0\Big).
\]
Un $\tilde P^{NL}$ que oscila a una frecuencia nueva actúa como fuente y radia luz a esa frecuencia (comparar con la ecuación de onda del tema 1).

\textbf{Simetría}: $\chi^{(2)}\ne0$ solo en cristales \textbf{no centrosimétricos}; $\chi^{(3)}$ existe en todos los materiales.

\subsection{Procesos de segundo orden (§1.2)}
\textbf{Un solo haz}: con $\tilde E=Ee^{-i\omega t}+\cc$, al elevar al cuadrado,
\[
\tilde P^{(2)}=2\epsilon_0\chi^{(2)}|E|^2+\big(\epsilon_0\chi^{(2)}E^2e^{-2i\omega t}+\cc\big)\qquad(1.2.2).
\]
El primer término es constante (rectificación óptica) y el segundo oscila a $2\omega$ (generación de segundo armónico).

\textbf{Dos haces}: con $\tilde E=E_1e^{-i\omega_1t}+E_2e^{-i\omega_2t}+\cc$, al expandir el cuadrado y escribir $\tilde P^{(2)}=\sum_n P(\omega_n)e^{-i\omega_nt}$:
\begin{clave}[title={Componentes de la polarización de segundo orden (1.2.7)}]
\renewcommand{\arraystretch}{1.5}
\begin{tabular}{@{}ll@{}}
$P(2\omega_1)=\epsilon_0\chi^{(2)}E_1^2$ & generación de segundo armónico (SHG)\\
$P(2\omega_2)=\epsilon_0\chi^{(2)}E_2^2$ & SHG\\
$P(\omega_1+\omega_2)=2\epsilon_0\chi^{(2)}E_1E_2$ & suma de frecuencias (SFG)\\
$P(\omega_1-\omega_2)=2\epsilon_0\chi^{(2)}E_1E_2^*$ & diferencia de frecuencias (DFG)\\
$P(0)=2\epsilon_0\chi^{(2)}\big(|E_1|^2+|E_2|^2\big)$ & rectificación óptica (OR)
\end{tabular}
\end{clave}
Los factores 2 en SFG y DFG aparecen porque los productos cruzados $E_1E_2$ y $E_2E_1$ salen dos veces al expandir el cuadrado. Si se usa la convención $\tilde E=\frac12(E'e^{-i\omega t}+\cc)$ (la del Saleh, $E'=2E$), todas las $P'$ llevan un factor $\tfrac12$ adicional (p.~7).

Puntos clave:
\begin{itemize}
\item Aunque la polarización tiene todas estas componentes, en la práctica solo una radia con eficiencia: aquella para la que se cumple el \textbf{ajuste de fase} (tema 5).
\item \textbf{Amplificación paramétrica} (DFG): por cada fotón de $\omega_3$ que se destruye se crean uno de $\omega_1$ y otro de $\omega_2$; la onda débil se amplifica. Sin semilla aparece la fluorescencia paramétrica; dentro de una cavidad se tiene el \textbf{OPO} (con ondas de bombeo, señal e \emph{idler}).
\item La rectificación óptica crea un campo eléctrico constante en el cristal, que no radia.
\item \textbf{Niveles virtuales}: en los diagramas de energía, los niveles intermedios no son autoestados del átomo. Un proceso es \textbf{paramétrico} si el estado cuántico final es igual al inicial (no hay transferencia de energía al medio y $\chi$ es real). Es \textbf{no paramétrico} si hay transferencia de población ($\chi$ compleja): absorción saturable $\alpha=\alpha_0/(1+I/I_s)$, absorción de dos fotones (sección eficaz $\sigma=\sigma^{(2)}I$, tasa $R=\sigma^{(2)}I^2/\hbar\omega$) y Raman estimulado.
\end{itemize}

\subsection{Procesos de tercer orden (§1.2.6–1.2.9)}
\textbf{Un solo haz} $\tilde E=\mathcal E\cos\omega t$: con la identidad $\cos^3\omega t=\frac14\cos3\omega t+\frac34\cos\omega t$,
\[
\tilde P^{(3)}=\tfrac14\epsilon_0\chi^{(3)}\mathcal E^3\cos3\omega t+\tfrac34\epsilon_0\chi^{(3)}\mathcal E^3\cos\omega t .
\]
El primer término es la \textbf{generación de tercer armónico} (THG). El segundo oscila a la misma frecuencia que el haz y se suma a la polarización lineal: cambia el índice de refracción (\textbf{efecto Kerr óptico}).

\begin{clave}[title={Índice no lineal}]
En la convención del Boyd, $P(\omega)=3\epsilon_0\chi^{(3)}|E|^2E$. Esto equivale a un índice que depende de la intensidad:
\[
n=n_0+\bar n_2\langle\tilde E^2\rangle=n_0+2\bar n_2|E|^2=n_0+n_2I,
\]
\[
\bar n_2=\frac{3\chi^{(3)}}{4n_0},\qquad \boxed{n_2=\frac{3\chi^{(3)}}{4n_0^2\epsilon_0c}}\quad(4.1.19);\qquad n_2[\mathrm{m^2/W}]=\frac{283\,\chi^{(3)}[\mathrm{m^2/V^2}]}{n_0^2}.
\]
\ojo{la ecuación (1.2.14b) imprime $3/(2n_0^2\epsilon_0c)$: es una errata del libro (sobra un factor 2); la correcta es (4.1.19).}

\textbf{Modulación de fase cruzada (XPM)}: un haz fuerte $E$ cambia el índice que ve una sonda débil $E'$: $P=6\epsilon_0\chi^{(3)}|E|^2E'$, así que $\bar n_2^{\rm cross}=2\bar n_2$ (el doble que la automodulación).
\end{clave}

\textbf{Tres frecuencias distintas}: el término $\tilde E^3$ produce 44 componentes de frecuencia (22 positivas y sus conjugadas). Regla para escribir cualquiera:
\[
P(\omega_\sigma)=D\,\epsilon_0\chi^{(3)}E_aE_bE_c ,
\]
donde $\omega_\sigma$ es la suma de las tres frecuencias (un campo conjugado aporta frecuencia negativa) y $D$ es el número de permutaciones distintas de los tres campos. Ejemplos:
\begin{itemize}
\item $P(3\omega_1)=\epsilon_0\chi^{(3)}E_1^3$ \quad ($D=1$, los tres campos iguales);
\item $P(2\omega_1\pm\omega_2)=3\epsilon_0\chi^{(3)}E_1^2E_2^{(*)}$ \quad ($D=3$, dos iguales);
\item $P(\omega_1+\omega_2\pm\omega_3)=6\epsilon_0\chi^{(3)}E_1E_2E_3^{(*)}$ \quad ($D=6$, todos distintos);
\item $P(\omega_1)=\epsilon_0\chi^{(3)}\big(3|E_1|^2+6|E_2|^2+6|E_3|^2\big)E_1$ \quad (automodulación más modulación cruzada).
\end{itemize}

\textbf{Autoenfoque}: si $n_2>0$ y el haz es gaussiano, el centro (más intenso) ve un índice mayor que los bordes. El medio actúa como una lente positiva y el haz se enfoca solo.

\subsection{Definición tensorial de la susceptibilidad (§1.3)}
En general $\mathbf P$ y $\mathbf E$ son vectores y $\chi^{(n)}$ es un tensor de rango $n+1$ que depende de las frecuencias. El campo se escribe $\tilde{\mathbf E}=\sum_n\mathbf E(\omega_n)e^{-i\omega_nt}$, con $\mathbf E(\omega_n)=\mathbf A_ne^{i\mathbf k_n\cdot\mathbf r}$.

\begin{clave}[title={Definición de $\chi^{(2)}_{ijk}$ y $\chi^{(3)}_{ijkl}$}]
\[
P_i(\omega_n+\omega_m)=\epsilon_0\sum_{jk}\sum_{(nm)}\chi^{(2)}_{ijk}(\omega_n+\omega_m,\omega_n,\omega_m)\,E_j(\omega_n)E_k(\omega_m)\qquad(1.3.12)
\]
La suma $(nm)$ recorre todas las combinaciones de frecuencias con suma $\omega_n+\omega_m$ fija. En la práctica se resume como
\[
P_i(\omega_n+\omega_m)=\epsilon_0D\sum_{jk}\chi^{(2)}_{ijk}E_jE_k,\qquad D=2\ (\text{SFG}),\quad D=1\ (\text{SHG}).
\]
Para tercer orden: $P_i(\omega_\sigma)=\epsilon_0D\sum_{jkl}\chi^{(3)}_{ijkl}(\omega_\sigma,\omega_o,\omega_n,\omega_m)E_jE_kE_l$, con $D=1,3,6$ (1.3.21).
\end{clave}
\ojo{en el límite $\omega_2\to\omega_1$, la SFG no se convierte en SHG de forma continua: $D$ pasa de 2 a 1.}

\subsection{Modelo clásico: oscilador anarmónico no centrosimétrico (§1.4.1)}
El modelo de Lorentz con una corrección cuadrática en la fuerza de restauración:
\[
\ddot x+2\gamma\dot x+\omega_0^2x+ax^2=-\frac{e\tilde E(t)}{m}\qquad(1.4.1),\qquad U(x)=\tfrac12m\omega_0^2x^2+\tfrac13max^3 .
\]
El término $x^3$ en el potencial lo hace asimétrico: el electrón se desplaza más fácil hacia un lado que hacia el otro. Esto solo es posible en un medio sin centro de inversión.

\textbf{Método de perturbaciones}: se reemplaza $\tilde E\to\lambda\tilde E$, se propone $x=\lambda x^{(1)}+\lambda^2x^{(2)}+\lambda^3x^{(3)}+\dots$ y se igualan potencias de $\lambda$. Con el operador $L=\partial_t^2+2\gamma\partial_t+\omega_0^2$ queda una ecuación lineal por orden:
\[
L[x^{(1)}]=-\frac{e\tilde E}{m},\qquad L[x^{(2)}]=-a\big(x^{(1)}\big)^2,\qquad L[x^{(3)}]=-2a\,x^{(1)}x^{(2)} .
\]
Cada orden se resuelve con la solución del anterior como fuente. Al final, $P=-Nex$.

\begin{clave}[title={Resultados del oscilador anarmónico}]
Denominador resonante: $D(\omega)=\omega_0^2-\omega^2-2i\omega\gamma$ (con $D(-\omega)=D(\omega)^*$).
\[
x^{(1)}(\omega_j)=-\frac{e}{m}\frac{E_j}{D(\omega_j)},\qquad \chi^{(1)}(\omega)=\frac{Ne^2/m}{\epsilon_0D(\omega)}\qquad(1.4.17)
\]
\[
x^{(2)}(2\omega_1)=\frac{-a(e/m)^2E_1^2}{D(2\omega_1)D^2(\omega_1)},\qquad x^{(2)}(\omega_1+\omega_2)=\frac{-2a(e/m)^2E_1E_2}{D(\omega_1+\omega_2)D(\omega_1)D(\omega_2)}
\]
\[
\chi^{(2)}(\omega_1+\omega_2,\omega_1,\omega_2)=\frac{N(e^3/m^2)\,a}{\epsilon_0D(\omega_1+\omega_2)D(\omega_1)D(\omega_2)}=\frac{\epsilon_0^2ma}{N^2e^3}\,\chi^{(1)}(\omega_1+\omega_2)\chi^{(1)}(\omega_1)\chi^{(1)}(\omega_2)\quad(1.4.24\text{–}25)
\]
\end{clave}
Otros procesos:
\[
\chi^{(2)}(2\omega,\omega,\omega)=\frac{N(e^3/m^2)a}{\epsilon_0D(2\omega)D^2(\omega)}\quad(\text{SHG});
\]
DFG, denominadores $D(\omega_1-\omega_2)D(\omega_1)D(-\omega_2)$; OR, $D(0)D(\omega)D(-\omega)$. \textbf{Regla general}: aparece un factor $1/D$ por cada frecuencia que interviene (las de entrada y la de salida). Por eso $\chi^{(2)}$ crece mucho si cualquiera de ellas se acerca a $\omega_0$.

\textbf{Regla de Miller}: el cociente
\[
\frac{\chi^{(2)}(\omega_3,\omega_1,\omega_2)}{\chi^{(1)}(\omega_3)\chi^{(1)}(\omega_1)\chi^{(1)}(\omega_2)}=\frac{\epsilon_0^2ma}{N^2e^3}
\]
es casi el mismo para todos los cristales. Es decir, los materiales con índice alto tienden a tener $\chi^{(2)}$ grande.

\textbf{Estimación de $\chi^{(2)}$}: si a un desplazamiento igual a la distancia interatómica $d$ los términos lineal y no lineal de la fuerza son comparables, $m\omega_0^2d=mad^2$, entonces $a=\omega_0^2/d$. Con $N=1/d^3$ y $D\to\omega_0^2$ (fuera de resonancia):
\[
\chi^{(2)}=\frac{e^3}{\epsilon_0m^2\omega_0^4d^4}\approx6.9\ \mathrm{pm/V}\qquad(\omega_0=10^{16}\ \mathrm{s^{-1}},\ d=3\ \text{Å}),
\]
del orden de lo medido en cristales reales.

\subsection{Oscilador centrosimétrico e isótropo (§1.4.3)}
En un medio con centro de inversión el potencial debe ser simétrico, así que la primera corrección es cúbica en la fuerza:
\[
\mathbf F=-m\omega_0^2\mathbf r+mb(\mathbf r\cdot\mathbf r)\mathbf r,\qquad U=\tfrac12m\omega_0^2x^2-\tfrac14mbx^4 .
\]
A segundo orden la ecuación para $\mathbf r^{(2)}$ queda homogénea (sin fuente), así que $\mathbf r^{(2)}=0$: no hay $\chi^{(2)}$, como exige la simetría. El primer término no lineal es de tercer orden:
\begin{clave}[title={$\chi^{(3)}$ del oscilador centrosimétrico}]
\[
\mathbf r^{(3)}(\omega_q)=-\sum_{(mnp)}\frac{b\,e^3\,[\mathbf E(\omega_m)\cdot\mathbf E(\omega_n)]\,\mathbf E(\omega_p)}{m^3D(\omega_q)D(\omega_m)D(\omega_n)D(\omega_p)},\qquad \mathbf P^{(3)}=-Ne\,\mathbf r^{(3)}
\]
\[
\chi^{(3)}_{ijkl}=\frac{Nbe^4\,[\delta_{ij}\delta_{kl}+\delta_{ik}\delta_{jl}+\delta_{il}\delta_{jk}]}{3\epsilon_0m^3D(\omega_q)D(\omega_m)D(\omega_n)D(\omega_p)}\qquad(1.4.52)\ (\text{forma simetrizada})
\]
\[
=\frac{bm\epsilon_0^3}{3N^3e^4}\,\chi^{(1)}(\omega_q)\chi^{(1)}(\omega_m)\chi^{(1)}(\omega_n)\chi^{(1)}(\omega_p)\,[\delta_{ij}\delta_{kl}+\delta_{ik}\delta_{jl}+\delta_{il}\delta_{jk}].
\]
Estimación con $b=\omega_0^2/d^2$: $\chi^{(3)}\approx\dfrac{e^4}{\epsilon_0m^3\omega_0^6d^5}\approx344$ pm²/V² ($\omega_0=7\times10^{15}$ s$^{-1}$, $d=3$ Å).
\end{clave}
Consecuencia de la estructura con deltas (Prob.~1.4): $\chi_{1122}=\chi_{1212}=\chi_{1221}=\frac13\chi_{1111}$, y todo elemento con un índice repetido un número impar de veces (p.~ej.\ $\chi_{1112}$) es cero.

\subsection{Propiedades de simetría de $\chi^{(2)}$ (§1.5.1–1.5.5)}
Las simetrías reducen mucho el número de elementos independientes del tensor.
\begin{itemize}
\item \textbf{Realidad de los campos}: como $\tilde P$ y $\tilde E$ son reales, $\chi_{ijk}(-\omega_n-\omega_m,-\omega_n,-\omega_m)=\chi_{ijk}(\omega_n+\omega_m,\omega_n,\omega_m)^*$.
\item \textbf{Permutación intrínseca} (siempre válida; es una convención): se pueden intercambiar a la vez el par índice–frecuencia de los dos campos de entrada, $\chi_{ijk}(\omega_n+\omega_m,\omega_n,\omega_m)=\chi_{ikj}(\omega_n+\omega_m,\omega_m,\omega_n)$.
\item \textbf{Medio sin pérdidas}: $\chi$ es real y vale la \textbf{permutación total}: se pueden permutar los tres pares índice–frecuencia, incluido el de salida (cambiando el signo de las frecuencias según corresponda):
$\chi_{ijk}(\omega_3=\omega_1+\omega_2)=\chi_{jki}(\omega_1=-\omega_2+\omega_3)=\chi_{kij}(\omega_2=\omega_3-\omega_1)$. Se demuestra a partir de la energía almacenada $U$, con $P_i=\partial U/\partial E_i^*$.
\item \textbf{Simetría de Kleinman}: si además las frecuencias están muy por debajo de las resonancias ($\omega\ll\omega_0$), $\chi$ casi no depende de la frecuencia y se pueden permutar los índices sin permutar las frecuencias:
$\chi_{ijk}=\chi_{jki}=\chi_{kij}=\chi_{ikj}=\chi_{jik}=\chi_{kji}$.
\item \textbf{Conteo de elementos}: 324 complejos en general $\to$ 81 $\to$ 27 (sin pérdidas) $\to$ 18 (SHG, matriz $d_{il}$) $\to$ 10 (Kleinman) $\to$ aún menos según la clase cristalina.
\end{itemize}

\subsection{Notación contraída y coeficiente efectivo $d_{\rm eff}$ (§1.5.6–1.5.7)}
Cuando vale la permutación intrínseca, $\chi_{ijk}$ es simétrico en $j,k$, y se define $d_{ijk}=\frac12\chi^{(2)}_{ijk}$. El par $jk$ se contrae en un solo índice $l$:
\[
11\to1,\quad22\to2,\quad33\to3,\quad23,32\to4,\quad13,31\to5,\quad12,21\to6 .
\]
Así $d_{ijk}$ se vuelve una matriz $3\times6$, $d_{il}$.

\begin{clave}[title={Polarización en notación contraída (1.5.26–1.5.29)}]
\textbf{SHG}:
\[
\begin{pmatrix}P_x\\P_y\\P_z\end{pmatrix}_{2\omega}=2\epsilon_0\,[d_{il}]\begin{pmatrix}E_x^2\\E_y^2\\E_z^2\\2E_yE_z\\2E_xE_z\\2E_xE_y\end{pmatrix}
\]
\textbf{SFG}: $\mathbf P(\omega_3)=4\epsilon_0[d_{il}]\big(E_{1x}E_{2x},\ E_{1y}E_{2y},\ E_{1z}E_{2z},\ E_{1y}E_{2z}+E_{1z}E_{2y},\ \ldots\big)^T$.

Para polarizaciones y dirección fijas, todo se reduce a un número, $d_{\rm eff}$:
\[
P(2\omega)=2\epsilon_0d_{\rm eff}E(\omega)^2,\qquad P(\omega_3)=4\epsilon_0d_{\rm eff}E(\omega_1)E(\omega_2).
\]
\end{clave}
\textbf{Kleinman en notación contraída}: $d_{21}=d_{16}$, $d_{24}=d_{32}$, $d_{31}=d_{15}$, $d_{14}=d_{25}=d_{36}$, $d_{23}=d_{34}$, $d_{12}=d_{26}$, $d_{13}=d_{35}$; quedan 10 elementos independientes.

\textbf{Clase 3m} (LiNbO$_3$, BBO):
\[
d=\begin{pmatrix}0&0&0&0&d_{31}&-d_{22}\\-d_{22}&d_{22}&0&d_{31}&0&0\\d_{31}&d_{31}&d_{33}&0&0&0\end{pmatrix}.
\]
Para un cristal uniaxial negativo:
tipo I ($o+o\to e$): $d_{\rm eff}=d_{31}\sin\theta-d_{22}\cos\theta\sin3\phi$; \ tipo II: $d_{\rm eff}=d_{22}\cos^2\theta\cos3\phi$ (1.5.30).

\textbf{Receta para calcular $d_{\rm eff}$}: (1) escribir los campos de entrada en los ejes del cristal según su polarización (o, e) y la dirección $(\theta,\phi)$; (2) calcular $\mathbf P$ con la matriz $[d]$; (3) proyectar $\mathbf P$ sobre la polarización de la onda de salida.

Otras clases: $\bar42m$ (KDP): $d_{14}=d_{25}$ y $d_{36}$. \ $\bar43m$ (GaAs) y 23: $d_{14}=d_{25}=d_{36}$. \ 432: todos los $d_{il}=0$.

\subsection{Simetría espacial del medio (§1.5.8–1.5.15)}
\textbf{Inversión implica $\chi^{(2)}=0$}: si el medio tiene centro de inversión, al invertir el campo se debe invertir la polarización: $-\tilde P=\epsilon_0\chi^{(2)}(-\tilde E)^2=\epsilon_0\chi^{(2)}\tilde E^2=\tilde P$. Solo se cumple si $\chi^{(2)}=0$.

De las 32 clases cristalinas, 11 son centrosimétricas y 21 no lo son. Las \textbf{polares} (con dipolo permanente) son 1, 2, 3, 4, 6, m, mm2, 3m, 4mm y 6mm.

Forma de la respuesta: un medio centrosimétrico solo genera armónicos impares y su respuesta promedio es nula; uno no centrosimétrico genera armónicos pares e impares y una componente DC.

Clasificación óptica lineal ($\chi^{(1)}$): cúbicos e isótropos son ópticamente isótropos; tetragonales, trigonales y hexagonales son uniaxiales; ortorrómbicos, monoclínicos y triclínicos son biaxiales.

Ejemplos: GaAs (zincblenda, $\bar43m$) tiene $d_{36}=370$ pm/V, pero como es cúbico no tiene birrefringencia y no se puede usar para ajuste de fase por birrefringencia. El Si (diamante, $m3m$) es centrosimétrico: $\chi^{(2)}_{\rm bulk}=0$.

\begin{center}
\renewcommand{\arraystretch}{1.25}
\begin{tabular}{@{}lll@{}}
\toprule
Cristal & Coeficiente & Valor (pm/V), Tabla 1.5.3\\\midrule
BBO & $d_{22}$ & 2.2\\
KDP & $d_{36}$ & 0.43\\
LiNbO$_3$ & $d_{31}$, $d_{33}$ & $-5.9$, \ $\approx-30$ (la tabla imprime “$d_{32}$”)\\
CdS & $d_{33}$ & 78\\
Cuarzo & $d_{11}$ & 0.3\\
AgGaSe$_2$ & $d_{36}$ & 33\\
\bottomrule
\end{tabular}
\end{center}

\textbf{$\chi^{(3)}$ en un medio isótropo}: 21 elementos no nulos y solo 3 independientes, ligados por $\chi_{1111}=\chi_{1122}+\chi_{1212}+\chi_{1221}$. Forma de Maker–Terhune (§4.2):
\[
\mathbf P^{(3)}(\omega)=\epsilon_0A(\mathbf E\cdot\mathbf E^*)\mathbf E+\tfrac12\epsilon_0B(\mathbf E\cdot\mathbf E)\mathbf E^*,\qquad A=6\chi_{1122},\quad B=6\chi_{1221}.
\]
El cociente $B/A$ revela el mecanismo físico: 1 para respuesta electrónica no resonante, 6 para orientación molecular, 0 para electrostricción. En cristales cúbicos (Si, GaAs) hay 4 elementos independientes: $\chi_{xxxx}$ ya no está fijado por los otros (anisotropía).

\subsection{Dominio del tiempo y relaciones de Kramers–Kronig (§1.6–1.7)}
Una respuesta no instantánea se describe con una función respuesta causal:
\[
\tilde P^{(1)}(t)=\epsilon_0\int_0^\infty R^{(1)}(\tau)\,\tilde E(t-\tau)\,d\tau,\qquad \chi^{(1)}(\omega)=\int_0^\infty R^{(1)}(\tau)e^{i\omega\tau}d\tau .
\]
A segundo orden, $\chi^{(2)}(\omega_\sigma;\omega_1,\omega_2)=\iint R^{(2)}(\tau_1,\tau_2)e^{i(\omega_1\tau_1+\omega_2\tau_2)}d\tau_1d\tau_2$. Transformada de Fourier en la convención del Boyd: $E(\omega)=\int\tilde E(t)e^{i\omega t}dt$. Una respuesta instantánea corresponde a $R\propto\delta(\tau)$, es decir, a $\chi$ independiente de la frecuencia.

Como $R(\tau)=0$ para $\tau<0$, $\chi(\omega)$ es analítica en el semiplano superior ($\mathrm{Im}\,\omega>0$) y cumple $\chi(\omega)=\frac{-i}{\pi}\mathrm{P}\!\int\frac{\chi(\omega')}{\omega'-\omega}d\omega'$. Separando partes real e imaginaria:
\begin{clave}[title={Kramers–Kronig en la convención del Boyd (1.7.9)}]
\[
\mathrm{Re}\,\chi^{(1)}(\omega)=\frac2\pi\,\mathrm P\!\int_0^\infty\frac{\omega'\,\mathrm{Im}\,\chi^{(1)}(\omega')}{\omega'^2-\omega^2}\,d\omega',\qquad
\mathrm{Im}\,\chi^{(1)}(\omega)=-\frac{2\omega}\pi\,\mathrm P\!\int_0^\infty\frac{\mathrm{Re}\,\chi^{(1)}(\omega')}{\omega'^2-\omega^2}\,d\omega'.
\]
\end{clave}
\textbf{Caso no lineal}: las relaciones valen si la respuesta es lineal en la frecuencia sobre la que se integra (todas las frecuencias de entrada se mueven juntas y son positivas). Valen para SHG $\chi^{(2)}(2\omega;\omega,\omega)$, THG y el índice inducido por un haz auxiliar $\chi^{(3)}(\omega;\omega,\omega_1,-\omega_1)$. \textbf{No} valen para el índice autoinducido $\chi^{(3)}(\omega;\omega,\omega,-\omega)$ ni para $\chi^{(2)}(\omega_1;\omega_3,-\omega_2)$.

\subsection{Conversión de unidades SI–gaussianas (Boyd, apéndice)}
Mucha literatura antigua usa unidades gaussianas (esu):
\begin{itemize}
\item $\chi^{(1)}_{SI}=4\pi\chi^{(1)}_{G}$;
\item $\chi^{(2)}[\mathrm{m/V}]=\dfrac{4\pi}{3\times10^4}\chi^{(2)}[\mathrm{esu}]=4.19\times10^{-4}\chi^{(2)}_G$;
\item $\chi^{(3)}[\mathrm{m^2/V^2}]=\dfrac{4\pi}{9\times10^8}\chi^{(3)}_G=1.40\times10^{-8}\chi^{(3)}_G$;
\item $d[\mathrm{esu}]=2.386\times10^{-9}\,d[\mathrm{pm/V}]$; \ $n_2[\mathrm{cm^2/W}]=0.0395\,\chi^{(3)}_G/n_0^2$.
\end{itemize}
Ejemplo: proustita, $1.3\times10^{-7}$ esu $\to5.4\times10^{-11}$ m/V.

\subsection{Cómo atacar un problema típico}
\begin{enumerate}
\item Escribir el campo en la convención del Boyd (cuidado: $\mathcal E=2|E|$).
\item Identificar la frecuencia de salida $\omega_\sigma$ y el factor de degeneración $D$.
\item Si hay tensores: llevar los campos a los ejes del cristal, aplicar $[d_{il}]$ o $\chi_{ijkl}$ con las simetrías de la clase, y proyectar sobre la polarización de salida.
\item Para modelos (oscilador): perturbación en $\lambda$, una ecuación por orden, $\mathbf P=-Ne\mathbf x$, e identificar $\chi$ con la definición (1.3.12) incluyendo $D$.
\item Revisar unidades y órdenes de magnitud (comparar con $E/E_{\rm at}$).
\end{enumerate}

\end{document}
